\chapter{Routing and navigation}
\label{ch:routing}

\begin{aside}
A TUI with multiple screens needs a way to move between them.
Whether that's tabs at the top, a back-stack, or URL-style
routes, the pattern is the same: which view is current?
\end{aside}

\section{The route signal}

A simple router:

\begin{lstlisting}
enum class Route { Home, Settings, About };
auto route = signal(Route::Home);

auto root = [&] {
    switch (ctx.get(route)) {
        case Route::Home:     return home();
        case Route::Settings: return settings();
        case Route::About:    return about();
    }
};
\end{lstlisting}

Setting the route signal swaps the visible view.

\section{The back stack}

For navigation history:

\begin{lstlisting}
class BackStack {
    std::vector<Route> stack;
public:
    void push(Route r) { stack.push_back(r); }
    void pop()        { if (!stack.empty()) stack.pop_back(); }
    Route current()    { return stack.back(); }
};
\end{lstlisting}

The current route is the top of the stack. Pop returns to
the previous view.

\section{Route parameters}

Some routes have parameters (which doc, which user):

\begin{lstlisting}
struct DocRoute { int doc_id; };
struct UserRoute { std::string username; };
using Route = std::variant<HomeRoute, DocRoute, UserRoute>;
\end{lstlisting}

The view function pattern-matches:

\begin{lstlisting}
auto view = std::visit([&](auto&& r) -> Element {
    using T = std::decay_t<decltype(r)>;
    if constexpr (std::is_same_v<T, HomeRoute>)
        return home();
    else if constexpr (std::is_same_v<T, DocRoute>)
        return doc_view(r.doc_id);
    // ...
}, ctx.get(route));
\end{lstlisting}

\section{Modal routes}

A modal (settings overlay over current view) is a route
that doesn't replace the current view but overlays it.
Tracked as a separate signal:

\begin{lstlisting}
auto modal = signal<std::optional<Modal>>(std::nullopt);
\end{lstlisting}

The view renders the modal if set, on top of the regular
view.

\section{Deep linking}

For apps that accept URL-like arguments, parse the route on
startup:

\begin{lstlisting}
int main(int argc, char** argv) {
    Route r = parse_route(argv[1]);
    ctx.set(route, r);
    run();
}
\end{lstlisting}

A \code{myapp settings} command-line jumps straight to the
settings view.

\section{Animated transitions}

Switching routes can animate (slide left, fade out + in).
Wrap the route view in an animator that crossfades on
change.

\section{Side effects on navigation}

Going to a route might trigger a fetch. Use an effect:

\begin{lstlisting}
ctx.effect([&] {
    auto r = ctx.get(route);
    if (auto* doc = std::get_if<DocRoute>(&r))
        fetch_document(doc->doc_id);
});
\end{lstlisting}

The effect re-runs on route change.

\section{Recap}

\begin{recap}
\begin{itemize}[leftmargin=*]
  \item Route signal selects the current view.
  \item Back stack for history.
  \item Variants for parameterised routes.
  \item Modals as a separate overlay route.
  \item Effects fire fetches on route change.
\end{itemize}
\end{recap}

\section*{Workshop}

\begin{enumerate}[leftmargin=*]
  \item Build a three-tab app with route switching.
  \item Add a back stack with Escape to pop.
  \item Crossfade between routes; tune the duration.
\end{enumerate}
